% !TeX root = ../main.tex
% ─────────────────────────────────────────────────────────────
%  EMPIRICAL EVALUATION — real numbers from experiments/run_eval.py
%  Source CSVs: experiments/results/{branching,gamelength_random,firstmove,
%               tournament,elo,timing,antirep}.csv  (see results/NOTES.md)
% ─────────────────────────────────────────────────────────────
\section{Empirical Evaluation}
\label{sec:eval}

This section quantifies three aspects of Cờ Trận: (i) the
breadth of the game tree via the \emph{branching factor} and game length, (ii) the relative
strength of the AI via a bot-vs-bot tournament, and (iii) real-time
performance. The goal is to confirm that the move space is \emph{broad and
non-degenerate} (branching factor), and that search depth \emph{actually
translates into playing strength} (greater depth beats lesser depth
consistently). Separation by depth is a sign, though not full proof,
that decisions in Cờ Trận have lasting consequences.

\subsection{Experimental Setup}
\label{sec:eval_setup}

All experiments run on an AMD Ryzen~7 6800H CPU (16 threads, 7.4\,GB RAM),
Python~3.12 on Linux (WSL2); the AI search runs single-threaded. The
measurement code lives in \texttt{experiments/} in the open-source
repository, with every random seed fixed for reproducibility. Because Cờ
Trận has no move limit and no draw rule (Section~\ref{sec:provenance}),
every game in the experiments of this section is played until one side
wins. The measurement code sets only a safety cap of $10^5$ plies against an
endless loop, and no game reached it. Each pairing
consists of 50 games; to balance the first-mover advantage, the two bots
alternate going first (each leads 25 games). The bots use the same
evaluation function from Equation~\eqref{eq:eval}, differing only in search
depth. Because each move at depth~6 costs several seconds on the pure Python
interpreter (see Section~\ref{sec:eval_timing}), the number of games is set
at 50 per pairing---enough for the depth-induced difference to show clearly
while staying within the compute budget.

\subsection{Branching Factor and State Space}
\label{sec:eval_branching}

Table~\ref{tab:branching} reports the average number of legal moves
(branching factor) by game phase (early/mid/late by thirds of game length).
It is measured over 60 complete games per mode, that is, about 20,000 sample
positions in the 2-player mode, 94,000 in the 3-player mode and 215,000 in
the 4-player mode. The sample
positions are generated by \emph{legal random play} (each turn picks a move
uniformly from the set of legal moves); the figures therefore reflect the
distribution of positions under random play, which may differ from the
distribution of positions that arise between two strong bots, but they give
an unbiased estimate of the breadth of the game tree. Under random play, a
game lasts a median of 304 plies in the 2-player mode, 1,326 plies in the
3-player mode and 3,259 plies in the 4-player mode. This does not measure
games between real people, but it is consistent with the recollection that
a four-player game used to last three to four hours
(Section~\ref{sec:provenance}).

\begin{table}[H]
  \centering
  \caption{Average branching factor (number of legal moves) by game phase
  and mode. Counting the choice of sandwich direction, the full number of
  \emph{actions} is only marginally higher (2-player $\approx 11.6$;
  3-player $\approx 18.9$; 4-player $\approx 21.1$). 60 random games per mode,
  each played to the end.}
  \label{tab:branching}
  \begin{tabular}{lccc}
    \toprule
    \textbf{Phase} & \textbf{2-player} & \textbf{3-player} & \textbf{4-player} \\
    \midrule
    Early game & 11.3 & 15.7 & 15.8 \\
    Mid game   & 11.1 & 19.0 & 21.7 \\
    Late game  & 11.1 & 20.5 & 22.8 \\
    \midrule
    \textbf{Game average} & 11.1 & 18.4 & 20.1 \\
    \bottomrule
  \end{tabular}
\end{table}

\noindent In the 2-player mode the average branching factor is $\approx 11$
and almost constant across game phases. This is a consequence of the
piece-conservation invariant: the piece count, and hence the number of
potential jumps, does not dwindle as in games with piece removal. In the 3-
and 4-player modes the branching factor \emph{grows} during the game, from
about 16 in the early game to 20--23 in the late game, and the 4-player mode
is the highest. We have not analysed the cause of this trend. All three
modes branch more than Reversi, which has about 10 moves per
turn~\citep{allis1994searching}. These values lie far below the analytic
upper bound $|\mathcal A(s)|\le 1408$ (Appendix~\ref{app:formal}),
confirming that the jump constraint tightens the actual move space;
positioning a game's complexity relative to solved games is standard
evaluation practice~\citep{herik2002games}, and Appendix~\ref{app:formal}
bounds the 2-player state space at about $3.7\times10^{17}$~\eqref{eq:statespace}. It must be stressed: the
branching factor measures the \emph{breadth} of the game tree, not directly
the \emph{depth} of tactics; it only shows that the move space is rich
enough and does not degenerate toward the endgame. The evidence for tactical
depth---that decisions truly matter---comes from the tournament in
Section~\ref{sec:eval_tournament}.

\subsection{Bot-vs-Bot Tournament}
\label{sec:eval_tournament}

To test whether the AI exploits search depth, we have bots of different
depths play adversarially (Table~\ref{tab:tournament}).

\begin{table}[H]
  \centering
  \caption{Bot-vs-bot results, 2-player mode (A is the deeper side; 50
  games per pairing, alternating first move 25/25). The 95\% confidence
  interval of the win-rate is computed by the Wilson method. The rules have
  no draws; every game in all three pairings ended with a winner.}
  \label{tab:tournament}
  \setlength{\tabcolsep}{4.5pt}
  \begin{tabular}{lccccc}
    \toprule
    \textbf{Pairing (depth)} & \textbf{Games} & \textbf{A wins}
      & \textbf{B wins} & \textbf{Win-rate A}
      & \textbf{95\% CI} \\
    \midrule
    d=6 vs d=2 & 50 & 50 & 0 & 100.0\,\% & [92.9,\,100.0] \\
    d=6 vs d=4 & 50 & 43 & 7 & 86.0\,\% & [73.8,\,93.0] \\
    d=4 vs d=2 & 50 & 45 & 5 & 90.0\,\% & [78.6,\,95.7] \\
    \bottomrule
  \end{tabular}
\end{table}

\paragraph{Elo Estimation.}
From the win-rate $p$ of the stronger side, the Elo difference is estimated
as $\Delta\mathrm{Elo} = -400\,\log_{10}(1/p - 1)$. Taking the depth-2 bot
as the 1000-Elo anchor, the d=4 vs d=2 pairing (90.0\%) gives
$\Delta\mathrm{Elo}\approx{+}382$, so depth~4 is $\approx 1382$ Elo, and the
d=6 vs d=4 pairing (86.0\%) gives ${+}315$. Chained together, depth~6 is
$\approx 1697$ Elo, about ${+}697$ above depth~2.\footnote{The d=6 vs d=2 pairing is a clean sweep (100\%),
which makes the logistic formula diverge ($p\to1$ sends
$\Delta\mathrm{Elo}\to\infty$), so it gives only a \emph{lower bound} rather
than a finite Elo value; we therefore estimate the depth-6 Elo via the
d=2\,$\to$\,d=4\,$\to$\,d=6 chain, which is more reliable.}

\paragraph{First-Mover Advantage.}
To probe a basic balance aspect of the 2-player mode, we have two bots of
\emph{equal depth} ($d{=}2$) self-play 100 games (fixed seed). The result is
notable: the \emph{second} mover (Yellow) wins 61 games and the \emph{first}
mover (Green) wins 39. Every game ended with a winner; the longest lasted
313 plies. The first mover's win rate is only $39.0\%$ (95\% Wilson
confidence interval $[30.0,\,48.8]$, lying entirely below $50\%$). Thus, for depth-2 bots,
\emph{moving second is advantageous} in a statistically significant way.
One hypothesis is that in a sandwich-flip game the second mover can
counterstrike after the opponent has revealed its position; we have not
tested this hypothesis. This
is an observation at \emph{one} fixed depth; the advantage may change with
depth and is a next step for the balance study.

\subsection{Real-Time Performance}
\label{sec:eval_timing}

Table~\ref{tab:timing} reports the per-move think time measured in
practice, for exactly the configurations in Table~\ref{tab:ai_params}.

\begin{table}[H]
  \centering
  \caption{Think time per move (seconds), measured over every move of 3
  complete self-play games per configuration (2-player: 919 moves;
  3-player: 2,029 moves; 4-player: 1,179 moves). The three configurations ran
  in parallel, one process per core; each search is still single-threaded,
  as in the game.}
  \label{tab:timing}
  \begin{tabular}{lccc}
    \toprule
    \textbf{Mode (depth)} & \textbf{Median} & \textbf{95th percentile}
      & \textbf{Max} \\
    \midrule
    2-player (d=6) & 3.07 & 5.80 & 9.51 \\
    3-player (d=4) & 1.11 & 2.35 & 4.47 \\
    4-player (d=3) & 0.17 & 0.57 & 1.40 \\
    \bottomrule
  \end{tabular}
\end{table}

\subsection{Discussion of Results}
\label{sec:eval_discussion}

Three observations emerge from the measurements above.

\textbf{(1) The move space is broad and non-degenerate.} A branching factor
of $\approx 11$ in the 2-player mode and $\approx 18$--$20$ in the 3- and
4-player modes (Table~\ref{tab:branching}) shows that each turn offers many
choices. In no mode does the branching shrink toward the end of the game: it
is stable in the 2-player mode and even grows in the 3- and 4-player modes,
unlike games with piece removal, where the branching dwindles toward the
endgame. The game therefore does not degenerate into a simple endgame.
Branching alone does not prove tactical depth; it establishes a necessary
condition (a non-trivial space), while the sufficient evidence comes from
observation~(2).

\textbf{(2) Search depth translates into playing strength (the evidence for
tactical depth).} The win-rate ranks monotonically with depth---d=6 beats
d=2 absolutely (100\%), d=4 beats d=2 (90\%), and d=6 beats d=4 (86\%)
(Table~\ref{tab:tournament})---every game ended with a winner, and all 95\%
confidence intervals lie clearly above $50\%$. That seeing farther brings a
systematic advantage indicates that decisions in Cờ Trận \emph{have
long-term consequences}---a sign of tactical depth---and that the evaluation
function combined with Minimax search is ``meaningful'', not random noise.
The Elo estimate above quantifies the gap: every two depth
levels correspond to $\approx 300$--$380$ Elo. Note that this is an
\emph{intrinsic} scale (relatively anchored among bots of the same engine),
not an Elo anchored to human players or an external opponent; it shows that
\emph{relative} strength grows with depth, not the absolute quality of the
AI.

\textbf{(3) Real-time performance.} The 4-player mode has a median of
$0.17$\,s per move, comfortable for interaction. The 3-player mode at
depth~4 takes a median of $1.11$\,s, at most $4.5$\,s. The 2-player mode at
depth~6 takes a median of $3.07$\,s and at most $9.5$\,s per move on the
single-threaded pure Python interpreter---beyond the threshold for instant
feedback. This is a strength--speed trade-off; it could be reduced by
adaptive depth reduction, time-limited search, or a compiled
implementation.

In summary, the experiments confirm both goals set at the start of the
section: the game \emph{does} have measurable tactical depth and the AI
\emph{plays reasonably}, with a clear time cost for the deepest search
configuration.
